两条等强度、理想极窄的平行狭缝间隔为 \(d>0\)，由远处空间扩展、各方向相互不相干但具有相同中心波长 \(\lambda>0\) 的光源照明。入射光的角度 \(\alpha\) 围绕零方向服从归一化高斯强度分布
\[
p(\alpha)=\frac1{\sqrt{2\pi}\sigma}\exp[-\alpha^2/(2\sigma^2)],
\]
其中 \(\sigma>0\) 足够小，观察屏的衍射角 \(\theta\) 也处于近轴范围。每个入射方向自身照射两缝是完全相干的，但不同 \(\alpha\) 的光强而非振幅相加。
(1) 对固定入射角 \(\alpha\)，写出两狭缝在观察方向 \(\theta\) 的远场干涉强度；(2) 对角分布积分求总光强与条纹可见度；(3) 定义条纹可见度衰减到 \(1/e\) 时的狭缝间距为横向相干长度，求其表达式；(4) 讨论点源极限 \(\sigma\to0^+\)、大缝距极限以及为何不应将不同入射角的复振幅直接相加。
